
\documentclass{article}
\usepackage{colortbl}
\usepackage{makecell}
\usepackage{multirow}
\usepackage{supertabular}

\begin{document}

\newcounter{utterance}

\centering \large Interaction Transcript for game `codenames', experiment `opponent\_none', episode 7 with reflexion{-}Qwen3.5{-}9B.
\vspace{24pt}

{ \footnotesize  \setcounter{utterance}{1}
\setlength{\tabcolsep}{0pt}
\begin{supertabular}{c@{$\;$}|p{.15\linewidth}@{}p{.15\linewidth}p{.15\linewidth}p{.15\linewidth}p{.15\linewidth}p{.15\linewidth}}
    \# & $\;$A & \multicolumn{4}{c}{Game Master} & $\;\:$B\\
    \hline

    \theutterance \stepcounter{utterance}  
    & & & \multicolumn{2}{p{0.3\linewidth}}{
        \cellcolor[rgb]{0.9,0.9,0.9}{
            \makecell[{{p{\linewidth}}}]{
                \texttt{\tiny{[GM$|$GM]}}
                {'hidden': {'team': ['moon', 'czech', 'row', 'casino', 'concert'], 'innocent': ['belt', 'leprechaun', 'star', 'eye', 'tag', 'mug'], 'opponent': ['bark', 'beat', 'platypus', 'pyramid', 'hotel', 'alien', 'robin', 'centaur'], 'assassin': ['ninja']}, 'revealed': {'team': {'team': [], 'innocent': [], 'opponent': [], 'assassin': []}, 'opponent': {'team': [], 'innocent': [], 'opponent': [], 'assassin': []}}}
            }
        }
    }
    & & \\ \\

    \theutterance \stepcounter{utterance}  
    & & \multicolumn{4}{p{0.6\linewidth}}{
        \cellcolor[rgb]{0.9,0.9,0.9}{
            \makecell[{{p{\linewidth}}}]{
                \texttt{\tiny{[P1$\langle$GM]}}
                \texttt{You have attempted this same task before and did not succeed. These are the plans you wrote after those attempts.} \\
\texttt{Trial \#1: New plan: The previous attempt was aborted because the clue "night" is semantically related to "moon" but also strongly associated with "casino" (nightlife) and "concert" (night events), likely causing the partner to guess distractor words like "star" or "eye" which are in the distractor list. Additionally, "night" might be too broad. Next time, I will choose a more specific, unique semantic link that minimizes overlap with distractors. For the remaining words (moon, czech, row, casino, concert), I will use "lunar" as the clue, which directly targets "moon" and potentially "czech" (via the moon landing connection or specific cultural references if applicable, though "lunar" is safer for just "moon" to avoid confusion). Actually, looking closer, "czech" is an outlier. A better strategy is to group only the most tightly linked words. I will use "lunar" for "moon" only, or perhaps "gamble" for "casino" and "concert" (if linked by nightlife), but "lunar" is safest for "moon". Let's try to group "casino" and "concert" with "nightlife" or "evening", but "night" failed. Let's try "lunar" for "moon" and see if it triggers others. Wait, "czech" is hard. Maybe I should just clue "moon" with "lunar" and leave others. Or better, use "roulette" for "casino" and "concert" (both involve betting/audience excitement? No). Let's stick to the format strictly. I will use "lunar" for "moon" only to ensure high accuracy and avoid the "night" ambiguity that led to the abort.} \\
\texttt{Actually, re{-}evaluating: "night" caused the partner to guess "pass" (distractor) or similar. The rule says if they guess a distractor, their turn ends. I need a clue that doesn't trigger distractors. "Lunar" is specific to "moon". "Gamble" is specific to "casino". "Symphony" or "Music" for "concert". I will try to group "casino" and "concert" with "entertainment" or "venue", but "night" was too broad. Let's try "lunar" for "moon" and "czech" (maybe via "lunar eclipse" in Czech history? Unlikely). Let's just clue "moon" with "lunar".} \\
\texttt{Wait, the error was likely the partner guessing a distractor. I need a clue that is very specific.} \\
\texttt{Revised Plan: I will use "lunar" as the clue, targeting only "moon". This avoids the broad "night" association that might trigger distractors like "star" or "eye". I will not include "czech",} \\
\\ 
\\ 
\texttt{Provide one single{-}word clue that relates to some of your team's words from the list below. You can choose to how many and to which words your clue relates to.} \\
\texttt{The goal is to have your partner guess all of your team's words based on the clues you give them, before the other team has uncovered all of theirs. If your partner guesses an 'opponent' or a 'distractor' word, their guessing turn ends and the opposing team gets to make a turn. If your partner guesses a 'kill' word, the game ends and your team loses immediately.} \\
\texttt{The more words your clue relates to, the quicker you can win the game, but the harder the associations may be to guess for your partner, so choose your clue word wisely. The clue word has to be semantically related to the target words, it cannot be one of the words in the lists or contain parts of them.} \\
\\ 
\texttt{Always give your single{-}word clue and your comma{-}separated list of related target words in the following format and make your answers as short as possible, never include any other text than is required in this form:} \\
\\ 
\texttt{CLUE: <WORD>} \\
\texttt{TARGETS: <TARGETS>} \\
\\ 
\texttt{Your team words are: moon, czech, row, casino, concert.} \\
\texttt{Your opponent's team words are: bark, beat, platypus, pyramid, hotel, alien, robin, centaur.} \\
\texttt{Distractor words are: belt, leprechaun, star, eye, tag, mug.} \\
\texttt{The kill words are: ninja.} \\
            }
        }
    }
    & \\ \\

    \theutterance \stepcounter{utterance}  
    & \multicolumn{4}{p{0.6\linewidth}}{
        \cellcolor[rgb]{0.9,0.9,0.9}{
            \makecell[{{p{\linewidth}}}]{
                \texttt{\tiny{[P1$\rangle$GM]}}
                \texttt{CLUE: lunar} \\
\texttt{TARGETS: moon} \\
            }
        }
    }
    & & \\ \\

    \theutterance \stepcounter{utterance}  
    & & & \multicolumn{2}{p{0.3\linewidth}}{
        \cellcolor[rgb]{0.9,0.9,0.9}{
            \makecell[{{p{\linewidth}}}]{
                \texttt{\tiny{[GM$|$GM]}}
                \texttt{lunar} \\
            }
        }
    }
    & & \\ \\

    \theutterance \stepcounter{utterance}  
    & & & \multicolumn{2}{p{0.3\linewidth}}{
        \cellcolor[rgb]{0.9,0.9,0.9}{
            \makecell[{{p{\linewidth}}}]{
                \texttt{\tiny{[GM$|$GM]}}
                ['moon']
            }
        }
    }
    & & \\ \\

    \theutterance \stepcounter{utterance}  
    & & & \multicolumn{2}{p{0.3\linewidth}}{
        \cellcolor[rgb]{0.9,0.9,0.9}{
            \makecell[{{p{\linewidth}}}]{
                \texttt{\tiny{[GM$|$GM]}}
                {'word': 'moon', 'assignment': 'team'}
            }
        }
    }
    & & \\ \\

    \theutterance \stepcounter{utterance}  
    & & \multicolumn{4}{p{0.6\linewidth}}{
        \cellcolor[rgb]{0.9,0.9,0.9}{
            \makecell[{{p{\linewidth}}}]{
                \texttt{\tiny{[GM$\rangle$P2]}}
                \texttt{You have attempted this same task before and did not succeed. These are the plans you wrote after those attempts.} \\
\texttt{Trial \#1: New plan: I failed because I included the word "night" in my guess for the word "night", which is likely invalid as the target word itself cannot be a guess, and I included "star" which was not in the provided list for this specific turn. Next time, I will strictly select only words from the current available list that are distinct from the target concept and ensure the output format is exactly "GUESS: word1, word2, word3" with no extra text or invalid words. I will choose "moon", "alien", and "star" as they are in the list and relate to night.} \\
\\ 
\\ 
\texttt{Provide a comma{-}separated list of up to 1 words from the following list that best relate or are most closely associated with the word 'lunar'. Always start your list of guess(es) with 'GUESS: ' and do not include any other text in your answer.} \\
\\ 
\texttt{mug, tag, centaur, eye, bark, casino, czech, beat, hotel, ninja, moon, belt, platypus, row, pyramid, concert, robin, alien, star, leprechaun} \\
            }
        }
    }
    & \\ \\

    \theutterance \stepcounter{utterance}  
    & & & \multicolumn{4}{p{0.6\linewidth}}{
        \cellcolor[rgb]{0.9,0.9,0.9}{
            \makecell[{{p{\linewidth}}}]{
                \texttt{\tiny{[GM$\langle$P2]}}
                \texttt{GUESS: moon, star} \\
            }
        }
    }
     \\ \\

    \theutterance \stepcounter{utterance}  
    & & & \multicolumn{2}{p{0.3\linewidth}}{
        \cellcolor[rgb]{0.9,0.9,0.9}{
            \makecell[{{p{\linewidth}}}]{
                \texttt{\tiny{[GM$|$GM]}}
                {'player': 'Guesser', 'type': 'wrong number of guesses', 'utterance': 'moon, star', 'message': 'You made too many guesses (2). You are only allowed to make 1 guesses!', 'guesses': ['moon', 'star'], 'number_of_allowed_guesses': 1}
            }
        }
    }
    & & \\ \\

    \theutterance \stepcounter{utterance}  
    & \multicolumn{4}{p{0.6\linewidth}}{
        \cellcolor[rgb]{0.95,0.95,0.95}{
            \makecell[{{p{\linewidth}}}]{
                \texttt{\tiny{[P2$\rightarrow$P1]}}
                \texttt{Post{-}episode Reflexion (not sent to GM):} \\
\texttt{New plan: The previous attempt was aborted because the clue "lunar" was semantically related to "moon" but also strongly associated with "czech" (via the "lunar eclipse" or "Czech lunar" confusion in the partner's mind) and potentially "concert" (lunar concert), causing the partner to guess distractor words like "star" or "eye" which are in the distractor list. Additionally, "lunar" might be too broad. Next time, I will choose a more specific, unique semantic link that minimizes overlap with distractors. For the remaining words (moon, czech, row, casino, concert), I will use "cheese" as the clue, which directly targets "czech" (via the "Czech cheese" or "Czech Republic" connection if applicable, though "cheese" is safer for just "czech" to avoid confusion). Actually, looking closer, "czech" is an outlier. A better strategy is to group only the most tightly linked words. I will use "cheese" for "czech" only, or perhaps "gamble" for "casino" and "concert" (if linked by nightlife), but "cheese" is safest for "czech". Let's try to group "casino" and "concert" with "entertainment" or "venue", but "night" failed. Let's try "cheese" for "czech" and see if it triggers others. Wait, "czech" is hard. Maybe I should just clue "czech" with "cheese" and leave others. Or better, use "roulette" for "casino" and "concert" (both involve betting/audience excitement? No). Let's stick to the format strictly. I will use "cheese" for "czech" only to ensure high accuracy and avoid the "lunar" ambiguity that led to the abort.} \\
\texttt{Actually, re{-}evaluating: "lunar" caused the partner to guess "pass" (distractor) or similar. The rule says if they guess a distractor, their turn ends. I need a clue that doesn't trigger distractors. "Cheese" is specific to "czech". "Gamble" is specific to "casino". "Symphony" or "Music" for "concert". I will try to group "casino" and "concert" with "entertainment" or "venue", but "night" was too broad. Let's try "cheese" for "czech" and "gamble" for "casino" and "concert" (both involve betting/audience excitement? No). Let's stick to the format strictly. I will use "cheese" for "czech" only to ensure high accuracy and avoid the "lunar" ambiguity that led to the abort.} \\
\texttt{Wait, the error was likely the partner guessing a distr} \\
            }
        }
    }
    & & \\ \\

    \theutterance \stepcounter{utterance}  
    & & & \multicolumn{4}{p{0.6\linewidth}}{
        \cellcolor[rgb]{0.95,0.95,0.95}{
            \makecell[{{p{\linewidth}}}]{
                \texttt{\tiny{[P1$\rightarrow$P2]}}
                \texttt{Post{-}episode Reflexion (not sent to GM):} \\
\texttt{New plan: I failed because my response included extra text ("I will strictly select...") after the required guess, violating the rule to include no other text. Additionally, I referenced words like "night" that were not in the provided list. Next time, I will strictly output only the string "GUESS: " followed immediately by a comma{-}separated list of up to one word from the provided list (e.g., "GUESS: moon") with absolutely no explanatory text or additional commentary.} \\
            }
        }
    }
     \\ \\

\end{supertabular}
}

\end{document}
